\documentclass[11pt,reqno]{amsart}

\usepackage[T1]{fontenc}
\usepackage[utf8]{inputenc}
\usepackage{amsmath,amssymb,amsthm,mathtools}
\usepackage{enumitem}
\usepackage{microtype}
\usepackage{xcolor}
\usepackage[colorlinks=true,
  linkcolor=blue!55!black,
  citecolor=green!40!black,
  urlcolor=blue!65!black]{hyperref}
\hypersetup{
  pdftitle={The Spherical Regge Scissors Theorem: A Self-Contained Proof through Weight-Two Motives},
  pdfauthor={}}

\input{paper/macros}

\newcommand{\Gm}{\mathbf G_m}
\newcommand{\DM}{\operatorname{DM}}
\newcommand{\Spec}{\operatorname{Spec}}
\newcommand{\Ext}{\operatorname{Ext}}
\newcommand{\gr}{\operatorname{gr}}
\newcommand{\Tr}{\operatorname{Tr}}

\allowdisplaybreaks

\title[The spherical Regge theorem]
  {The Spherical Regge Scissors Theorem:\\
   A Self-Contained Proof through Weight-Two Motives}
\author{}
\date{Research draft, 27 September 2026}

\begin{document}

\begin{abstract}
We prove that every nondegenerate spherical tetrahedron is scissors
congruent to each elementary Regge mate.  This text contains only the proof
of that statement and the background needed to read it.  In particular, we
derive the monomial Regge action on angle phases, rationalize the orientation
double cover of the Gram-parameter space, explain the particular relative
quadric motive used in the proof, compute its three weight pieces and its
coproduct, and reduce the sole primitive obstruction to
\(H^1(U,\Q(2))=0\).  We then spell out the specialization to the given
tetrahedron and the final passage from Goncharov's quadric invariant to a
finite spherical equidecomposition.  General motivic theory and alternative
chain constructions are not developed beyond the precise standard results
used in this argument.
\end{abstract}

\maketitle
\tableofcontents

\input{paper/sections/01_statement}
\input{paper/sections/02_phase_charts}
\input{paper/sections/03_motivic_dictionary}
\input{paper/sections/04_relative_rigidity}
\input{paper/sections/05_specialization_scissors}

\bibliographystyle{alpha}
\bibliography{paper/biblio}

\end{document}
